\documentclass[11pt]{article}
\usepackage{mathblatt}


\blattfuss{Wachstum}{Lösungen}
\begin{document}
\noindent{\large\bfseries Wachstum \textperiodcentered{} Lösungen}\par\medskip
\begin{pfloesung}{}
\lz{1.}{$55\,€$}{neuer Wert: $50 \cdot 1{,}1$ $\Rightarrow$ 55}{}
\end{pfloesung}
\begin{pfloesung}{}
\lz{2.}{$60\,\mathrm{kg}$}{neuer Wert: $80 \cdot 0{,}75$ $\Rightarrow$ 60}{}
\end{pfloesung}
\begin{pfloesung}{}
\lz{3.}{$11$}{y}{}
\end{pfloesung}
\begin{pfloesung}{}
\lz{4.}{$A(1 | 2)$ und $B(3 | 4)$ eingetragen}{}{}
\end{pfloesung}
\begin{pfloesung}{}
\lz{5.}{$6$}{}{}
\end{pfloesung}
\begin{pfloesung}{}
\lz{6.}{$12$}{}{}
\end{pfloesung}
\begin{pfloesung}{}
\lz{7.}{$30\,€$}{}{}
\end{pfloesung}
\begin{pfloesung}{}
\lz{8.}{$30\,\%$}{}{}
\end{pfloesung}
\begin{pfloesung}{}
\lz{9.}{$1{,}4$}{}{}
\end{pfloesung}
\begin{pfloesung}{}
\lz{10.}{$3$}{$\tfrac{84}{28}$ $\Rightarrow$ 3}{}
\end{pfloesung}
\begin{pfloesung}{}
\lz{11.}{$32$}{}{}
\end{pfloesung}
\begin{pfloesung}{}
\lz{12.}{$3{,}375$}{}{}
\end{pfloesung}
\begin{pfloesung}{}
\lz{13.}{der Quotient}{}{}
\end{pfloesung}
\begin{pfloesung}{}
\lz{14.}{$439{,}23$}{$300 \cdot 1{,}1^4$ $\Rightarrow$ 439{,}23}{}
\end{pfloesung}
\begin{pfloesung}{}
\lz{15.}{der Startwert steht im Jahr $0$, er wird nicht noch einmal mit $1{,}1$ multipliziert}{300}{}
\end{pfloesung}
\begin{pfloesung}{}
\lz{16.}{$300$}{alter Wert: 375 : 1{,}25 $\Rightarrow$ 300}{}
\end{pfloesung}
\begin{pfloesung}{}
\lz{17.}{$287{,}5$}{neuer Wert: $250 \cdot 1{,}15$ $\Rightarrow$ 287{,}5}{}
\end{pfloesung}
\begin{pfloesung}{}
\lz{18.}{$343$}{700; 490}{}
\end{pfloesung}
\begin{pfloesung}{}
\lz{19.}{$14\,739$}{14\,739\,€: 24\,000; 20\,400; 17\,340}{}
\end{pfloesung}
\begin{pfloesung}{}
\lz{20.}{$1\,331$}{Jahr: $3,$ denn $1\,000 \cdot 1{,}1^3$}{}
\end{pfloesung}
\begin{pfloesung}{}
\lz{21.}{$484$}{484: $440 \cdot 1{,}1$ $\Rightarrow$ 484; Probe: $484 \cdot 1{,}1$ $\Rightarrow$ 532{,}4}{}
\end{pfloesung}
\begin{pfloesung}{{\small\color{mbgrau}P10 2020 · 4a}}
\lz{22.}{2019: 54; 2022: 57,3 · 1,03 = 59,019 $\approx$ 59,0}{neuer Wert: 57,3 · 1,03 $\Rightarrow$ 59,019 $\approx$ 59,0}{}
\end{pfloesung}
\begin{pfloesung}{{\small\color{mbgrau}P10 2016 · 4a}}
\lz{23.}{4,45 mg; 4 h}{neuer Wert: 5,00 · 0,89 $\Rightarrow$ 4,45}{}
\end{pfloesung}
\begin{pfloesung}{{\small\color{mbgrau}P10 2017 · 7a}}
\lz{24.}{756,9 $\approx$ 757 hPa; $\approx$ 328 hPa}{neuer Wert: 870 · 0,87 $\Rightarrow$ 756,9 $\approx$ 757}{}
\end{pfloesung}
\begin{pfloesung}{{\small\color{mbgrau}P10 2019 · 7a}}
\lz{25.}{10 kg; 10,816 kg; $\approx$ 12,167 kg}{Wachstumsfaktor: 1,04; neuer Wert: 10,4 · 1,04 $\Rightarrow$ 10,816}{}
\end{pfloesung}
\begin{pfloesung}{}
\lz{26.}{$350$}{}{}
\end{pfloesung}
\begin{pfloesung}{{\small\color{mbgrau}P10 2025 · 7a}}
\lz{27.}{Punkte (0|80), (1|112), (2|157), (3|220), (4|308) eingetragen, zunehmend steiler}{Skala wählen}{}
\end{pfloesung}
\begin{pfloesung}{{\small\color{mbgrau}P10 2021 · 6b}}
\lz{28.}{x-Achse 1 Kästchen = 5 min, y-Achse 1 Kästchen = 2 cm; fallende Strecke von (0; 40) bis (80; 24)}{80 min $\mathrel{\widehat{=}}$ 16 Kästchen; 40 cm $\mathrel{\widehat{=}}$ 20 Kästchen}{}
\end{pfloesung}
\begin{pfloesung}{}
\lz{29.}{$192$}{50; 98; Ablesegenauigkeit: $\pm 5$}{}
\end{pfloesung}
\begin{pfloesung}{}
\lz{30.}{Punkte $(0 | 100)$, $(1 | 120)$, $(2 | 144)$, $(3 | 173)$, $(4 | 207)$}{}{}
\end{pfloesung}
\begin{pfloesung}{}
\lz{31.}{$125$ liegt genau in der Mitte zwischen $100$ und $150$}{Punkte: (0 | 125), (1 | 150), (2 | 180), (3 | 216)}{}
\end{pfloesung}
\begin{pfloesung}{}
\lz{32.}{glatte, nach rechts steiler werdende Kurve durch $(0 | 100)$ bis $(4 | 207)$}{}{}
\end{pfloesung}
\begin{pfloesung}{}
\lz{33.}{Punkte eingetragen}{waagerecht: 1 Jahr je Kästchen, senkrecht 100\,€ je Kästchen; 531 passt in 800}{}
\end{pfloesung}
\begin{pfloesung}{{\small\color{mbgrau}P10 2016 · 4b}}
\lz{34.}{x-Achse Zeit in h, 1 h je Teilstrich; y-Achse Masse in mg, 1 mg je Teilstrich; sieben Punkte, fallende Kurve}{Skala x: 0 bis 8, y 0 bis 7}{}
\end{pfloesung}
\begin{pfloesung}{{\small\color{mbgrau}P10 2017 · 7b}}
\lz{35.}{x-Achse Höhe in km, 2 Kästchen je km; y-Achse Luftdruck in hPa, 1 Kästchen je 50 hPa; sieben Punkte, fallende Kurve}{x 0 bis 10 km, y 0 bis 1000 hPa}{}
\end{pfloesung}
\begin{pfloesung}{}
\lz{36.}{der Faktor ist größer als eins}{mehr}{}
\end{pfloesung}
\begin{pfloesung}{}
\lz{37.}{$2016$ ist der Schritt null}{5}{}
\end{pfloesung}
\begin{pfloesung}{}
\lz{38.}{$\approx 1{,}2167$}{$1{,}04^5$ $\Rightarrow$ 1{,}22}{}
\end{pfloesung}
\begin{pfloesung}{}
\lz{39.}{$1{,}4$}{}{}
\end{pfloesung}
\begin{pfloesung}{}
\lz{40.}{$1\,536$}{2\,400; 1\,920}{}
\end{pfloesung}
\begin{pfloesung}{}
\lz{41.}{$40$ und $1{,}3$}{Anfangswert und Faktor}{}
\end{pfloesung}
\begin{pfloesung}{}
\lz{42.}{Abnahme um $10\,\%$ heißt Faktor $0{,}9$}{$2\,000 \cdot 0{,}9^t$}{}
\end{pfloesung}
\begin{pfloesung}{}
\lz{43.}{$600 \cdot 0{,}8^t$}{m(t)}{}
\end{pfloesung}
\begin{pfloesung}{}
\lz{44.}{$1\,536\,€$}{$3\,000 \cdot 0{,}8^3$ $\Rightarrow$ 1536}{}
\end{pfloesung}
\begin{pfloesung}{}
\lz{45.}{$0{,}9^2$, Faktor $0{,}9$}{10\,\%: 810 : 1\,000 $\Rightarrow$ 0{,}81}{}
\end{pfloesung}
\begin{pfloesung}{}
\lz{46.}{$40 \cdot 1{,}3^t$}{A(t)}{}
\end{pfloesung}
\begin{pfloesung}{}
\lz{47.}{$1{,}2$, also $20\,\%$ Zunahme}{420 : 350 $\Rightarrow$ 1{,}2}{}
\end{pfloesung}
\begin{pfloesung}{}
\lz{48.}{Zunahme um $15\,\%$}{}{}
\end{pfloesung}
\begin{pfloesung}{}
\lz{49.}{Startwert $80$ und Quotient $1{,}25$}{$80 \cdot 1{,}25^t$}{}
\end{pfloesung}
\begin{pfloesung}{}
\lz{50.}{$0{,}75$}{}{}
\end{pfloesung}
\begin{pfloesung}{}
\lz{51.}{$1{,}05$: Faktor, jedes Jahr $5\,\%$ mehr}{250: 250: Mitglieder zu Beginn}{}
\end{pfloesung}
\begin{pfloesung}{{\small\color{mbgrau}P10 2017 · 7c}}
\lz{52.}{y = 1000 · $0{,}87^{x}$}{}{}
\end{pfloesung}
\begin{pfloesung}{{\small\color{mbgrau}P10 2018 · 2a}}
\lz{53.}{600 000 € · 1,08 = 648 000 €}{48 000 € : 600 000 € $\Rightarrow$ 8 \%}{}
\end{pfloesung}
\begin{pfloesung}{{\small\color{mbgrau}P10 2019 · 7b}}
\lz{54.}{10: Anfangsmasse in kg; 1,04: Wachstumsfaktor, +4 \% je Woche; x: Zeit in Wochen}{}{}
\end{pfloesung}
\begin{pfloesung}{}
\lz{55.}{der Bestand wächst jedes Jahr mit dem Faktor $1{,}2$, also um $20\,\%$}{1{,}2; 1{,}2; 1{,}2}{}
\end{pfloesung}
\begin{pfloesung}{}
\lz{56.}{$\approx 3\,546$}{$6: 2\,500 \cdot 1{,}06^6$ $\Rightarrow$ 0}{}
\end{pfloesung}
\begin{pfloesung}{}
\lz{57.}{die Behauptung stimmt nicht, es fehlen knapp $40\,€$}{$2\,000 \cdot 1{,}04^{10}$ $\Rightarrow$ $\approx$ 2\,960{,}49\,€}{}
\end{pfloesung}
\begin{pfloesung}{{\small\color{mbgrau}P10 2017 · 7d}}
\lz{58.}{ja, 1000 · 0,87$^{4}$ $\approx$ 573 hPa}{0,87$^{4}$ $\approx$ 0,573}{}
\end{pfloesung}
\begin{pfloesung}{{\small\color{mbgrau}P10 2016 · 4e}}
\lz{59.}{$\approx$ 0,31 mg}{}{}
\end{pfloesung}
\begin{pfloesung}{{\small\color{mbgrau}P10 2020 · 4c}}
\lz{60.}{nein, Studie 1 54 · 1,03$^{14}$ $\approx$ 81,7 gegen Studie 2 = 75,6 (Tausend Liter)}{neuer Wert: 54 · 1,03$^{14}$ $\Rightarrow$ $\approx$ 81,7; neuer Wert: 54 · 1,4 $\Rightarrow$ 75,6}{}
\end{pfloesung}
\begin{pfloesung}{{\small\color{mbgrau}P10 2025 · 7b}}
\lz{61.}{112 : 80 = 1,4 (alle Quotienten $\approx$ 1,4); N(t) = 80 · $1{,}4^{t}$; wahr, N(24) $\approx$ 257 136 > 250 000}{$1,4^24$ $\Rightarrow$ $\approx$ 3214,2}{}
\end{pfloesung}
\begin{pfloesung}{}
\lz{62.}{$100$}{}{}
\end{pfloesung}
\begin{pfloesung}{}
\lz{63.}{fallende Kurve, die immer flacher wird}{f}{}
\end{pfloesung}
\begin{pfloesung}{}
\lz{64.}{$f$ passt: Start bei $100$, gekrümmt. $g$ passt nicht: Gerade. $h$ passt nicht: Start im Ursprung, die Fläche ist aber schon zu Beginn $100$}{}{}
\end{pfloesung}
\begin{pfloesung}{}
\lz{65.}{Start bei $300$ und Gerade, weil jede Minute derselbe Betrag dazukommt}{g}{}
\end{pfloesung}
\begin{pfloesung}{}
\lz{66.}{$f$ passt: Start bei $45$, Kurve. $g$ passt nicht: Gerade, also jede Woche gleich viel. $h$ passt nicht: Start bei $0$, das Kalb wiegt aber $45\,\mathrm{kg}$}{}{}
\end{pfloesung}
\begin{pfloesung}{}
\lz{67.}{$b$ ist die Gerade}{}{}
\end{pfloesung}
\begin{pfloesung}{}
\lz{68.}{eine Gerade hat immer denselben Zuwachs}{Gleicher Prozentsatz heißt gleicher Faktor: die Zuwächse werden von Jahr zu Jahr größer und der Graph steiler}{}
\end{pfloesung}
\begin{pfloesung}{{\small\color{mbgrau}P10 2019 · 7c}}
\lz{69.}{Graph 1 nein (linear); Graph 2 ja (Start 10 kg, Zuwächse wachsen); Graph 3 nein (Start bei 0 kg)}{Startwert: 10 kg; gleicher Prozentsatz $\Rightarrow$ wachsende Zuwächse}{}
\end{pfloesung}
\begin{pfloesung}{\textbf{Prüfstein} \textperiodcentered{} P10 2026 FOR · 7}
\lz{a)}{2026: 650,00 €; 2029: 687,76 €}{Wachstumsfaktor: 1,019; neuer Wert: 674,93 · 1,019 $\Rightarrow$ 687{,}75}{}
\lz{b)}{B; A: Miete beginnt bei 0 € statt bei 650 €; C: linear, jedes Jahr gleicher Betrag statt gleicher Prozentsatz}{Startwert: 650 €; konstanter Prozentsatz $\Rightarrow$ wachsende Zuwächse}{}
\lz{c)}{M(t) = 650 · $1{,}019^{t}$; M(14) $\approx$ 845,96 € (rund 846 €)}{Unterschied: 2040 $-$ 2026 $\Rightarrow$ 14; $1,019^14$ $\Rightarrow$ $\approx$ 1,3015}{}
\end{pfloesung}
\end{document}